\Needspace{14\baselineskip}

\CNsection{1}{Circuit Switching Fundamentals}

\CNchained\CNsubsection{}{Concept: Dedicated Path}

In circuit switching, a \textbf{dedicated communication path} is established between two endpoints for the entire duration of the call. This path is reserved exclusively — no other user can use those resources until the call is released.

\Needspace{17\baselineskip}

\CNsubsection{}{How It Works in GSM}

{\def\LTcaptype{none} % do not increment counter
\Needspace{13\baselineskip}
\begin{longtable}[c]{@{}
  >{\raggedright\arraybackslash\hspace{0pt}}m{(\linewidth - 2\tabcolsep) * \real{0.4286}}
  >{\raggedright\arraybackslash\hspace{0pt}}m{(\linewidth - 2\tabcolsep) * \real{0.5714}}@{}}
\toprule\noalign{}
\rowcolor{CNink}\CNth{Step} & \CNth{Action} \\
\CNheadrulesep
\endhead
\bottomrule\noalign{}
\endlastfoot
1 & Caller initiates a call \\
2 & Network allocates a dedicated timeslot on the air interface (TDMA) \\
3 & A fixed path is reserved through BTS \ensuremath{\to} BSC \ensuremath{\to} MSC \ensuremath{\to} trunk to destination \\
4 & Resources remain allocated for the entire call duration \\
5 & Upon hangup, all resources are released \\
\end{longtable}
}

\CNsubsection{}{Resource Reservation}

\begin{itemize}
\tightlist
\item
  \textbf{Air Interface}: A dedicated TDMA timeslot (one of 8 per carrier frequency) is assigned
\item
  \textbf{Abis Interface} (BTS\ensuremath{\leftrightarrow}BSC): A 16~kbps or 64~kbps channel on E1/\allowbreak{}T1 links
\item
  \textbf{A Interface} (BSC\ensuremath{\leftrightarrow}MSC): A 64~kbps timeslot on E1/\allowbreak{}T1 links
\item
  \textbf{Trunk} (MSC\ensuremath{\leftrightarrow}PSTN/\allowbreak{}other MSC): A 64~kbps circuit on PCM links
\end{itemize}

\CNsubsection{}{Key Characteristics}

\begin{itemize}
\tightlist
\item
  \textbf{Connection-oriented}: Path established before data flows
\item
  \textbf{Guaranteed bandwidth}: Fixed 64~kbps per direction once connected
\item
  \textbf{Constant delay}: Predictable latency (\textasciitilde{}150ms end-to-end target)
\item
  \textbf{Inefficient for bursty data}: Resources wasted during silence periods
\end{itemize}

\Needspace{14\baselineskip}

\CNsection{2}{SS7 Signaling Concepts}

\CNchained\CNsubsection{}{What is SS7?}

\textbf{Signaling System 7 (SS7)} is the out-of-band signaling protocol stack used by telecom networks to:

\begin{itemize}
\tightlist
\item
  Set up and tear down calls
\item
  Manage subscriber databases
\item
  Enable roaming and handovers
\item
  Provide supplementary services (call forwarding, SMS, etc.)
\end{itemize}

SS7 signaling is \textbf{separate} from the voice path — this is called \textbf{Common Channel Signaling (CCS)}.

\CNkeepfig{m02-ss7-stack}{3}

\CNsubsection{}{SS7 Protocol Stack (Simplified)}

\CNfigure{m02-ss7-stack}{Simplified SS7 protocol stack}

\Needspace{11\baselineskip}

\CNsubsection{}{Key Protocols}

\CNchained\CNsubsubsection{ISUP (ISDN User Part)}

\begin{itemize}
\tightlist
\item
  \textbf{Purpose}: Call setup, management, and release between exchanges (MSCs)
\item
  \textbf{Messages}: IAM (Initial Address Message), ACM (Address Complete), ANM (Answer), REL (Release), RLC (Release Complete)
\item
  \textbf{Used for}: Establishing the voice circuit between switches
\end{itemize}

\CNsubsubsection{MAP (Mobile Application Part)}

\begin{itemize}
\tightlist
\item
  \textbf{Purpose}: Communication between GSM network elements (MSC, HLR, VLR, AuC)
\item
  \textbf{Operations}: Location updates, authentication data retrieval, subscriber info queries, SMS routing
\item
  \textbf{Used for}: Mobility management, roaming, subscriber data exchange
\end{itemize}

\CNsubsubsection{TCAP (Transaction Capabilities Application Part)}

\begin{itemize}
\tightlist
\item
  \textbf{Purpose}: Provides a framework for MAP and other applications to exchange structured data
\item
  \textbf{Model}: Request-response transactions (invoke, result, error)
\item
  \textbf{Used for}: Carrying MAP operations over the SS7 network
\end{itemize}

\Needspace{13\baselineskip}

\CNsubsection{}{Signaling Points}

{\def\LTcaptype{none} % do not increment counter
\Needspace{9\baselineskip}
\begin{longtable}[c]{@{}
  >{\raggedright\arraybackslash\hspace{0pt}}m{(\linewidth - 2\tabcolsep) * \real{0.6000}}
  >{\raggedright\arraybackslash\hspace{0pt}}m{(\linewidth - 2\tabcolsep) * \real{0.4000}}@{}}
\toprule\noalign{}
\rowcolor{CNink}\CNth{Element} & \CNth{Role} \\
\CNheadrulesep
\endhead
\bottomrule\noalign{}
\endlastfoot
\textbf{SSP} (Service Switching Point) & MSC/\allowbreak{}GMSC — originates/\allowbreak{}terminates signaling \\
\textbf{STP} (Signal Transfer Point) & Routes signaling messages between SSPs \\
\textbf{SCP} (Service Control Point) & Database nodes (HLR, AuC) \\
\end{longtable}
}

\Needspace{14\baselineskip}

\CNsection{3}{Call Setup Process}

\CNchained\CNsubsection{3.1}{Mobile-Originated Call (MOC)}

A Mobile-Originated Call is initiated by the mobile subscriber.

\CNsubsubsection{Sequence of Events}

\begin{enumerate}
\def\labelenumi{\arabic{enumi}.}
\tightlist
\item
  \textbf{UE dials number} \ensuremath{\to} sends CHANNEL REQUEST on RACH
\item
  \textbf{BTS} forwards to BSC
\item
  \textbf{BSC} assigns a dedicated channel (SDCCH) \ensuremath{\to} IMMEDIATE ASSIGNMENT
\item
  \textbf{UE} sends CM SERVICE REQUEST (service type = MOC) to MSC via BSC
\item
  \textbf{MSC} authenticates the subscriber (see Section 6)
\item
  \textbf{MSC} activates ciphering
\item
  \textbf{UE} sends SETUP message with called party number
\item
  \textbf{MSC} sends CALL PROCEEDING to UE
\item
  \textbf{MSC} analyzes digits, determines routing
\item
  \textbf{MSC} sends IAM (ISUP) toward GMSC/\allowbreak{}PSTN
\item
  \textbf{GMSC/\allowbreak{}PSTN} returns ACM (Address Complete) \ensuremath{\to} ringback
\item
  \textbf{MSC} sends ALERTING to UE
\item
  \textbf{Called party answers} \ensuremath{\to} ANM (Answer Message) returned
\item
  \textbf{MSC} sends CONNECT to UE
\item
  \textbf{BSC} assigns a TCH (Traffic Channel) for voice
\item
  \textbf{Voice path established} — conversation begins
\end{enumerate}

\CNkeepfig{m02-moc-1}{2}

\CNsubsubsection{MOC Sequence Diagram}

\CNfigure{m02-moc-1}{Mobile-originated call (MOC) set-up\CNpart{1}{2}}

\CNfigure{m02-moc-2}{Mobile-originated call (MOC) set-up\CNpart{2}{2}}

\CNsubsection{3.2}{Mobile-Terminated Call (MTC)}

A Mobile-Terminated Call is a call arriving for a mobile subscriber.

\CNsubsubsection{Sequence of Events}

\begin{enumerate}
\def\labelenumi{\arabic{enumi}.}
\tightlist
\item
  \textbf{PSTN/\allowbreak{}calling party} sends IAM to GMSC with called MSISDN
\item
  \textbf{GMSC} queries HLR: ``Where is this subscriber?'' (MAP: Send Routing Info)
\item
  \textbf{HLR} returns MSRN (Mobile Station Roaming Number) pointing to serving MSC
\item
  \textbf{GMSC} routes call to serving MSC using MSRN (IAM)
\item
  \textbf{MSC} queries VLR for subscriber state and location area
\item
  \textbf{MSC} pages the UE in the Location Area
\item
  \textbf{UE} responds to paging on RACH
\item
  \textbf{BSC} assigns SDCCH \ensuremath{\to} authentication \& ciphering
\item
  \textbf{MSC} sends SETUP to UE
\item
  \textbf{UE} sends CALL CONFIRMED
\item
  \textbf{UE} alerts user (ringtone)
\item
  \textbf{User answers} \ensuremath{\to} UE sends CONNECT
\item
  \textbf{MSC} sends ANM back toward GMSC/\allowbreak{}PSTN
\item
  \textbf{Traffic channel assigned} — voice path active
\end{enumerate}

\CNkeepfig{m02-mtc-1}{2}

\CNsubsubsection{MTC Sequence Diagram}

\CNfigure{m02-mtc-1}{Mobile-terminated call (MTC) set-up\CNpart{1}{2}}

\CNfigure{m02-mtc-2}{Mobile-terminated call (MTC) set-up\CNpart{2}{2}}

\Needspace{14\baselineskip}

\CNsection{4}{Call Routing}

\CNchained\CNsubsection{}{How MSC/\allowbreak{}GMSC Route Calls Using HLR}

The HLR is the \textbf{master database} of all subscribers. It knows which MSC/\allowbreak{}VLR currently serves each subscriber.

\CNsubsection{}{Routing a Mobile-Terminated Call}

\tcbset{CNcodemode/.style={unbreakable}}\def\CNcodelang{}

\begin{CNplaincode}
\begin{Verbatim}
Step 1: PSTN delivers call to GMSC (based on MSISDN number plan)
Step 2: GMSC sends MAP "Send Routing Information" to HLR
Step 3: HLR identifies the serving MSC/VLR for that subscriber
Step 4: HLR asks serving VLR to allocate an MSRN (temporary routing number)
Step 5: VLR returns MSRN to HLR
Step 6: HLR returns MSRN to GMSC
Step 7: GMSC uses MSRN to route the ISUP IAM to the correct MSC
\end{Verbatim}
\end{CNplaincode}

\CNsubsection{}{MSRN (Mobile Station Roaming Number)}

\begin{itemize}
\tightlist
\item
  Temporary number in E.164 format
\item
  Looks like a regular phone number in the serving network’s numbering plan
\item
  Allows GMSC to route the call using standard ISUP without knowing GSM internals
\item
  Allocated per call, released after routing
\end{itemize}

\CNsubsection{}{Routing a Mobile-Originated Call}

For MOC, routing is simpler:

\begin{enumerate}
\def\labelenumi{\arabic{enumi}.}
\tightlist
\item
  MSC receives the dialed digits from the UE
\item
  MSC performs digit analysis (number plan, prefixes, supplementary services)
\item
  MSC routes the call via ISUP toward the appropriate trunk/\allowbreak{}GMSC/\allowbreak{}PSTN gateway
\item
  If calling another mobile on same network \ensuremath{\to} MSC may query HLR for MSRN
\end{enumerate}

\Needspace{14\baselineskip}

\CNsection{5}{Location Update Procedure}

\CNchained\CNsubsection{}{Why Location Updates Are Needed}

The network must know \textbf{where a subscriber is} to deliver calls and messages. Without location tracking:

\begin{itemize}
\tightlist
\item
  The network would have to page every cell in the entire country
\item
  This would overwhelm the paging channel capacity
\item
  Call setup time would be unacceptable
\end{itemize}

Location updates allow the network to \textbf{narrow paging to a specific Location Area} (group of cells).

\Needspace{15\baselineskip}

\CNsubsection{}{Types of Location Updates}

{\def\LTcaptype{none} % do not increment counter
\Needspace{11\baselineskip}
\begin{longtable}[c]{@{}
  >{\raggedright\arraybackslash\hspace{0pt}}m{(\linewidth - 2\tabcolsep) * \real{0.4000}}
  >{\raggedright\arraybackslash\hspace{0pt}}m{(\linewidth - 2\tabcolsep) * \real{0.6000}}@{}}
\toprule\noalign{}
\rowcolor{CNink}\CNth{Type} & \CNth{Trigger} \\
\CNheadrulesep
\endhead
\bottomrule\noalign{}
\endlastfoot
\textbf{IMSI Attach} & Phone powers on — informs network subscriber is reachable \\
\textbf{Normal LU} & UE moves to a new Location Area (detects new LAI on BCCH) \\
\textbf{Periodic LU} & Timer-based update to confirm UE is still reachable \\
\textbf{IMSI Detach} & Phone powers off gracefully — informs network \\
\end{longtable}
}

\CNkeepfig{m02-location-update}{3}

\CNsubsection{}{Location Update Sequence Diagram}

\CNfigure{m02-location-update}{Location update with a change of MSC/\allowbreak{}VLR}

\CNsubsection{}{IMSI Attach}

When a phone powers on:

\begin{enumerate}
\def\labelenumi{\arabic{enumi}.}
\tightlist
\item
  UE reads BCCH to get LAI and network information
\item
  UE sends LOCATION UPDATING REQUEST with attach flag
\item
  Same procedure as normal LU follows
\item
  VLR marks subscriber as ``attached'' (reachable)
\end{enumerate}

\CNsubsection{}{Periodic Location Update}

\begin{itemize}
\tightlist
\item
  Network broadcasts a timer value (T3212) on BCCH
\item
  UE must send a LU before the timer expires, even if it hasn’t moved
\item
  If the timer expires without an LU, VLR marks subscriber as ``implicitly detached''
\item
  Purpose: detect phones that powered off without IMSI Detach (battery died, etc.)
\end{itemize}

\Needspace{14\baselineskip}

\CNsection{6}{Authentication Procedure}

\CNchained\CNsubsection{}{Purpose}

Authentication ensures that the subscriber is who they claim to be, preventing:

\begin{itemize}
\tightlist
\item
  Unauthorized network access
\item
  Cloned SIM fraud
\item
  Eavesdropping (when combined with ciphering)
\end{itemize}

\Needspace{15\baselineskip}

\CNsubsection{}{The Authentication Triplet}

The AuC (Authentication Center, co-located with HLR) generates \textbf{authentication triplets}:

{\def\LTcaptype{none} % do not increment counter
\Needspace{9\baselineskip}
\begin{longtable}[c]{@{}ccc@{}}
\toprule\noalign{}
\rowcolor{CNink}\CNth{Element} & \CNth{Size} & \CNth{Purpose} \\
\CNheadrulesep
\endhead
\bottomrule\noalign{}
\endlastfoot
\textbf{RAND} & 128~bits & Random challenge sent to UE \\
\textbf{SRES} & 32~bits & Expected response (computed from RAND + Ki) \\
\textbf{Kc} & 64~bits & Cipher key for encrypting the air interface \\
\end{longtable}
}

\CNkeepfig{m02-triplets}{3}

\CNsubsection{}{How Triplets Are Generated}

\CNfigure{m02-triplets}{Generation of an authentication triplet (RAND, SRES, Kc)}

\begin{itemize}
\tightlist
\item
  \textbf{Ki} never leaves the SIM card or the AuC
\item
  The same computation happens independently on both sides
\item
  If SRES from UE matches SRES from triplet \ensuremath{\to} subscriber is authenticated
\end{itemize}

\CNkeepfig{m02-authentication}{3}

\CNsubsection{}{Authentication Sequence Diagram}

\CNfigure{m02-authentication}{GSM authentication and ciphering}

\CNsubsection{}{Key Points}

\begin{itemize}
\tightlist
\item
  VLR stores \textbf{multiple triplets} (typically 5) to avoid querying HLR for every transaction
\item
  Authentication is performed at: call setup, location update, SMS, supplementary service activation
\item
  GSM authentication is \textbf{one-way} — network authenticates the UE, but UE does not authenticate the network (vulnerability fixed in 3G/\allowbreak{}UMTS)
\end{itemize}

\Needspace{14\baselineskip}

\Needspace{20\baselineskip}

\CNsection{7}{Roaming}

\CNchained\CNsubsection{}{Concepts}

{\def\LTcaptype{none} % do not increment counter
\Needspace{11\baselineskip}
\begin{longtable}[c]{@{}
  >{\raggedright\arraybackslash\hspace{0pt}}m{(\linewidth - 2\tabcolsep) * \real{0.3333}}
  >{\raggedright\arraybackslash\hspace{0pt}}m{(\linewidth - 2\tabcolsep) * \real{0.6667}}@{}}
\toprule\noalign{}
\rowcolor{CNink}\CNth{Term} & \CNth{Definition} \\
\CNheadrulesep
\endhead
\bottomrule\noalign{}
\endlastfoot
\textbf{HPLMN} & Home PLMN — the network where the subscriber has their subscription \\
\textbf{VPLMN} & Visited PLMN — a foreign network the subscriber is using while roaming \\
\textbf{Roaming Agreement} & Commercial + technical agreement between two operators \\
\textbf{GRX/\allowbreak{}IPX} & Interconnect network for signaling between PLMNs \\
\end{longtable}
}

\CNsubsection{}{How Roaming Works}

\begin{enumerate}
\def\labelenumi{\arabic{enumi}.}
\tightlist
\item
  \textbf{UE powers on in foreign country} \ensuremath{\to} selects VPLMN (based on SIM’s preferred PLMN list)
\item
  \textbf{UE performs Location Update} to VPLMN’s MSC/\allowbreak{}VLR
\item
  \textbf{VPLMN’s VLR} contacts \textbf{HPLMN’s HLR} via SS7/\allowbreak{}MAP (through international signaling links)
\item
  \textbf{HLR} authenticates subscriber (sends triplets to VPLMN)
\item
  \textbf{HLR} sends subscriber profile to VPLMN’s VLR
\item
  \textbf{HLR} records the VPLMN’s VLR address as the subscriber’s current location
\item
  \textbf{Subscriber can now make/\allowbreak{}receive calls} through VPLMN
\end{enumerate}

\CNkeepfig{m02-roaming}{3}

\CNsubsection{}{Signaling Path for Roaming}

\CNfigure{m02-roaming}{Signalling path for roaming: visited MSC/\allowbreak{}VLR to home HLR/\allowbreak{}AuC over SS7/\allowbreak{}MAP}

\CNsubsection{}{Incoming Call to a Roaming Subscriber}

\begin{enumerate}
\def\labelenumi{\arabic{enumi}.}
\tightlist
\item
  Call arrives at HPLMN’s GMSC
\item
  GMSC queries HLR \ensuremath{\to} learns subscriber is in VPLMN
\item
  HLR requests MSRN from VPLMN’s VLR
\item
  VPLMN’s VLR returns MSRN
\item
  HLR gives MSRN to GMSC
\item
  GMSC routes call to VPLMN’s MSC using MSRN (via international trunks)
\item
  VPLMN’s MSC pages and connects the subscriber
\end{enumerate}

\CNsubsection{}{Outgoing Call from Roaming Subscriber}

\begin{enumerate}
\def\labelenumi{\arabic{enumi}.}
\tightlist
\item
  UE initiates call in VPLMN
\item
  VPLMN’s MSC handles the call locally
\item
  Call is routed from VPLMN to destination (may transit through HPLMN or go direct)
\item
  CAMEL/\allowbreak{}CAP protocol may be used for prepaid control from HPLMN
\end{enumerate}

\Needspace{14\baselineskip}

\CNsection{8}{Call Release}

\CNchained\CNsubsection{}{Normal Call Release (Mobile-Initiated)}

\begin{enumerate}
\def\labelenumi{\arabic{enumi}.}
\tightlist
\item
  \textbf{UE} sends DISCONNECT to MSC
\item
  \textbf{MSC} sends REL (Release) via ISUP toward the far end
\item
  \textbf{Far end} responds with RLC (Release Complete)
\item
  \textbf{MSC} sends RELEASE to UE
\item
  \textbf{UE} sends RELEASE COMPLETE
\item
  \textbf{BSC} deallocates the TCH (traffic channel)
\item
  \textbf{All resources freed}: timeslots, trunks, circuit references
\end{enumerate}

\CNsubsection{}{Network-Initiated Release}

Can be triggered by:

\begin{itemize}
\tightlist
\item
  Far-end hangup
\item
  Radio link failure (UE out of coverage)
\item
  BSC/\allowbreak{}MSC timer expiry
\item
  Administrative disconnect
\end{itemize}

\CNkeepfig{m02-isup-release}{3}

\CNsubsection{}{ISUP Release Messages}

\CNfigure{m02-isup-release}{ISUP call release}

\Needspace{16\baselineskip}

\CNsection{9}{Why Circuit Switching?}

\CNchained\CNsubsection{}{Why Circuit Switching Was Used for Voice}

\CNchained\CNsubsubsection{Voice Requirements}

\begin{itemize}
\tightlist
\item
  \textbf{Constant bit rate}: 13~kbps (full rate) or 6.5~kbps (half rate) — continuous stream
\item
  \textbf{Low and predictable delay}: \textless{} 150ms one-way for acceptable conversation quality
\item
  \textbf{No jitter tolerance}: Human ear is very sensitive to variable delay
\item
  \textbf{Symmetric traffic}: Both parties speak roughly equally
\end{itemize}

\Needspace{16\baselineskip}

\CNsubsubsection{Why Circuit Switching Fits Voice}

{\def\LTcaptype{none} % do not increment counter
\Needspace{13\baselineskip}
\begin{longtable}[c]{@{}cc@{}}
\toprule\noalign{}
\rowcolor{CNink}\CNth{Requirement} & \CNth{How CS Addresses It} \\
\CNheadrulesep
\endhead
\bottomrule\noalign{}
\endlastfoot
Low delay & Dedicated path = no queuing delay \\
No jitter & Fixed timeslot allocation = constant delay \\
Guaranteed bandwidth & Reserved resources = no congestion drops \\
Simple processing & Once path is set up, switching is trivial \\
Real-time & No store-and-forward delays \\
\end{longtable}
}

\CNsubsubsection{Historical Context (1980s-90s)}

\begin{itemize}
\tightlist
\item
  Processing power was expensive — simple switching was preferred
\item
  Memory was expensive — no buffering/\allowbreak{}packetization
\item
  Existing PSTN was circuit-switched — natural integration
\item
  Voice was the \textbf{only} service expected from mobile phones
\end{itemize}

\Needspace{11\baselineskip}

\Needspace{20\baselineskip}

\CNsubsection{}{Why Circuit Switching Became a Limitation}

\CNchained\CNsubsubsection{The Problem with Data}

{\def\LTcaptype{none} % do not increment counter
\Needspace{14\baselineskip}
\begin{longtable}[c]{@{}
  >{\raggedright\arraybackslash\hspace{0pt}}m{(\linewidth - 2\tabcolsep) * \real{0.3500}}
  >{\raggedright\arraybackslash\hspace{0pt}}m{(\linewidth - 2\tabcolsep) * \real{0.6500}}@{}}
\toprule\noalign{}
\rowcolor{CNink}\CNth{Issue} & \CNth{Explanation} \\
\CNheadrulesep
\endhead
\bottomrule\noalign{}
\endlastfoot
\textbf{Resource waste} & A 64~kbps circuit is held even during silence (50-60\% of a voice call) \\
\textbf{No statistical multiplexing} & Cannot share resources among multiple users dynamically \\
\textbf{Fixed bandwidth} & Cannot burst above 64~kbps even when network is idle \\
\textbf{Call setup delay} & 3-7 seconds to establish a circuit — unacceptable for web browsing \\
\textbf{Per-minute charging} & Inefficient for checking email (30 seconds of data, charged 1 minute) \\
\textbf{Scalability} & Each data session consumes a full voice circuit \\
\end{longtable}
}

\CNkeepfig{m02-traffic-patterns}{2}

\CNsubsubsection{Data Traffic Characteristics vs. Voice}

\CNfigure{m02-traffic-patterns}{Occupancy of a reserved circuit: continuous voice versus bursty data}

Circuit switching \textbf{cannot efficiently handle bursty traffic} because resources are reserved but unused most of the time.

\CNkeepfig{m02-evolution}{2}

\CNsubsubsection{The Evolution Path}

\CNfigure{m02-evolution}{From circuit switching to fully packet-switched networks}

\CNsubsection{}{Summary}

Circuit switching was the \textbf{right choice} for voice in the 1990s — it provided the quality guarantees that voice communication demands with the technology available at that time. However, as mobile usage shifted from voice-only to data-centric services, the inefficiency of holding dedicated circuits for bursty, asymmetric data traffic made packet switching the inevitable successor.

\Needspace{28\baselineskip}

\CNsection{}{Key Terms Summary}

{\def\LTcaptype{none} % do not increment counter
\Needspace{22\baselineskip}
\begin{longtable}[c]{@{}
  >{\raggedright\arraybackslash\hspace{0pt}}m{(\linewidth - 2\tabcolsep) * \real{0.3333}}
  >{\raggedright\arraybackslash\hspace{0pt}}m{(\linewidth - 2\tabcolsep) * \real{0.6667}}@{}}
\toprule\noalign{}
\rowcolor{CNink}\CNth{Term} & \CNth{Definition} \\
\CNheadrulesep
\endhead
\bottomrule\noalign{}
\endlastfoot
\textbf{MOC} & Mobile-Originated Call \\
\textbf{MTC} & Mobile-Terminated Call \\
\textbf{MSRN} & Mobile Station Roaming Number — temporary number for call routing \\
\textbf{IAM} & Initial Address Message — ISUP message to initiate a call \\
\textbf{ACM} & Address Complete Message — called party is being alerted \\
\textbf{ANM} & Answer Message — called party has answered \\
\textbf{REL/\allowbreak{}RLC} & Release /\allowbreak{} Release Complete — call teardown \\
\textbf{LAI} & Location Area Identity — identifies a group of cells \\
\textbf{TMSI} & Temporary Mobile Subscriber Identity — privacy protection \\
\textbf{Ki} & Secret key shared between SIM and AuC \\
\textbf{SRES} & Signed Response — authentication proof \\
\textbf{Kc} & Cipher key derived during authentication \\
\textbf{VPLMN} & Visited PLMN — foreign network during roaming \\
\textbf{HPLMN} & Home PLMN — subscriber’s home network \\
\end{longtable}
}

\emph{Next Module: {Module 3 — GPRS and Packet Switching}}
